Language models depend on the evidence supplied to their prediction and generation processes. Tables can omit information needed for semantic annotation, classification labels can evolve under limited supervision, and answers can require relationships distributed across textual and visual sources. This thesis investigates how explicit knowledge structures can organize and select evidence for these settings.

The thesis develops four methods for structural augmentation. KGLink integrates knowledge graphs with a pre-trained language model to enrich column representations and bridge differences between graph types and annotation labels. GORAG supports dynamic few-shot classification through connected label retrieval and an online graph index that evolves with the task. AGRAG selects connected textual evidence by combining query-dependent node influence, semantic edge costs, and hybrid passage retrieval. MGRAG constructs query-specific multimodal graphs, matches query relations to semantic paths, and uses graph structure to guide the reuse and recomputation of cached model states.

Experiments across table annotation, dynamic classification, textual generation, and multimodal question answering show improved effectiveness over established baselines and efficiency gains in graph indexing and context processing. Ablation studies verify the contributions of the structural components. A related study of hyper-relational knowledge graphs examines information preservation. These studies demonstrate how explicit relationships improve evidence representation, retrieval, and computation for language models.
